% Source: https://artofproblemsolving.com/community/c3h1763286_harder_than_it_looks
% Stable identity: aops-topic-1763286-q01
Label the vertices $V_0,\dots,V_5$ in order, with the leash tied at $V_0$; $V_3$ is the opposite vertex. A point $P$ outside the barn is reachable when the shortest path from $V_0$ to $P$ that goes around the barn has length at most $600$. Such a path runs along the boundary, clockwise or counterclockwise, to some vertex $V_k$ and then straight to $P$.

\textbf{Halving the region.} The region is symmetric about the line $V_0V_3$ (dashed). For a point $P$ on the side of $V_1,V_2$, going round that way is never longer, so the half of the region on this side consists of
\begin{itemize}
\item half of the sector at $V_0$: the leash sweeps the exterior angle $240^\circ$ with radius $600$, and half of it is $\dfrac{1}{3}\pi\cdot600^2=120000\pi$;
\item the sector at $V_1$: the leash has $480$ left and sweeps the exterior angle $60^\circ$: $\dfrac{1}{6}\pi\cdot480^2=38400\pi$;
\item the part, on this side of the line $V_0V_3$, of the sector at $V_2$ with radius $360$ and angle $60^\circ$, bounded by the line $V_1V_2$ extended and the side $V_2V_3$ extended (green). Beyond the line $V_0V_3$ the other direction round takes over. A path wrapping round $V_3$ is never needed: going the other way round, via $V_4$, is at least as short there.
\end{itemize}

\begin{asy}
usepackage("amsmath");
size(8cm);
defaultpen(fontsize(10pt));
real s=1.2, L=6, R=L-2s;
pair[] V; for(int k=0;k<6;++k) V[k]=s*dir(-90+60*k);
path hex=V[0]--V[1]--V[2]--V[3]--V[4]--V[5]--cycle;
real ph=degrees(acos(sqrt(3)/6));
filldraw(arc(V[0],L,150,390)--V[0]--cycle,palegray,gray);
filldraw(V[1]--arc(V[1],L-s,30,90)--cycle,paleblue,gray);
filldraw(V[5]--arc(V[5],L-s,90,150)--cycle,paleblue,gray);
filldraw(V[2]--arc(V[2],R,90,180-ph)--V[3]--cycle,palegreen,gray);
filldraw(V[4]--V[3]--arc(V[4],R,ph,90)--cycle,palegreen,gray);
pair X=V[2]+R*dir(180-ph);
filldraw(hex,white,black+linewidth(1));
draw(V[0]--(0,-L-s),dashed);draw(V[3]--X,dashed);
draw(V[0]--V[0]+L*dir(-60),dashed+gray);draw(V[1]--V[1]+(L-s)*dir(45),dashed+gray);
dot(V[0]);dot(X);
label("$V_0$",V[0],dir(210));label("$V_1$",V[1],dir(-45));label("$V_5$",V[5],dir(225));
label("$V_2$",V[2],E);label("$V_4$",V[4],W);label("$V_3$",V[3],dir(35));label("$X$",X,N);
label("$600$",V[0]+3.5*dir(-60),dir(30));
label("$480$",V[1]+2.6*dir(45),dir(-45));label("$360$",V[2]+(0,2.4),E);
label("\mbox{}",(min(currentpicture,user=true)+max(currentpicture,user=true))/2,black); // keeps the text after the figure black when LaTeX typesets the labels inline
\end{asy}

\textbf{The clipped sector.} Put $V_3$ at the origin with the line $V_3V_0$ as the vertical axis. Then $V_2$ is at horizontal distance $120\cos30^\circ=60\sqrt3$ from the axis, and the sector at $V_2$ consists of the directions from straight up (along $V_1V_2$ extended) to the direction of $V_2V_3$, which makes $30^\circ$ with the horizontal. Measure the angle $\theta$ from the upward vertical towards $V_3$, so $\theta$ runs from $0$ to $60^\circ$. In direction $\theta$ the axis is at distance $\dfrac{60\sqrt3}{\sin\theta}$ from $V_2$, and the arc of radius $360$ meets the axis at $X$, where $\sin\theta_X=\dfrac{60\sqrt3}{360}=\dfrac{\sqrt3}{6}$. The part of the sector on our side is the sector minus the piece beyond the axis:
\[
21600\pi-\int_{\theta_X}^{\dfrac{\pi}{3}}\dfrac{1}{2}\left(360^2-\dfrac{10800}{\sin^2\theta}\right)d\theta
=21600\pi-64800\left(\dfrac{\pi}{3}-\theta_X\right)+5400\left(\cot\theta_X-\cot\dfrac{\pi}{3}\right).
\]
Here $\cot\theta_X=\sqrt{11}$, $\cot\dfrac{\pi}{3}=\dfrac{1}{\sqrt3}$, and $\dfrac{\pi}{3}-\theta_X=\arctan\sqrt{11}-\dfrac{\pi}{6}$ (because $\dfrac{\pi}{2}-\theta_X$ is the angle whose tangent is $\sqrt{11}$).

\textbf{Total.} The whole area is twice the half:
\[
2\left[120000\pi+38400\pi+21600\pi-64800\left(\arctan\sqrt{11}-\dfrac{\pi}{6}\right)+5400\sqrt{11}-1800\sqrt3\right],
\]
that is,
\[
\boxed{381600\pi-129600\arctan\sqrt{11}+10800\sqrt{11}-3600\sqrt3}\approx1.0628\times10^{6}.
\]
(Adding the full sectors $240000\pi+2\cdot38400\pi+2\cdot21600\pi=360000\pi\approx1.131\times10^6$ would count the overlap of the two sectors behind $V_3$ twice.)
